\scanpage{3}
\mspage{2}
context of a workshop defining a ``country''.

The preceding work has made it possible to bring out the most convenient point of view for the algebraic formalism. The fundamental notion does seem to me to be that of \uline{figure}, and not that of ``elementary figure'' or ``multistratum'', or even that of ``locus'' (or even, the notion of ``punctual locus'' or ``point'', which has not yet been discussed\ldots). One of the reasons for this is the simplicity of the algebraic formalism,
% slanted
\margin{The other reason is that, in the case of a ``map $f$'' from a ``region'' into another, the inverse image of a ``figure'' (\unsure{when|when} \unsure{one} can give a \unsure{geometric} meaning \unsure{to} \unsure{this} \unsure{word|term}) \deleted{\unsure{will be|becomes|having to}} tends \unsure{not} to \unsure{still} be a figure, but the inverse image of a multistratum is a multistratum.}%
 where it is pleasant to be able to manipulate loci and multistrata as being particular figures. When a figure \inserted{$F$} is interpreted \inserted{(by the set of its ``strata'')} as a subset of the set $\unsure{M|\mathcal M}$ \inserted{(ordered by $\leq$)} of multistrata, we shall write $\widetilde F$ (and not $F$) for this subset. When it is interpreted as a \inserted{set-theoretic} figure in $M$, we write it $\mathrm{Mulomb}(F)$ (``large multishadow''), its support is $\mathrm{Omb}(F)$ (the ``large shadow'' of $F$) -- the latter, unlike $\widetilde F$ and $\mathrm{Mulomb}(F)$, practically never makes it possible to reconstruct $F$, \deleted{\unsure{it|the}} but morally, its ``support'' only
